\documentclass[11pt,addpoints]{exam}
\usepackage{amsfonts,amssymb,amsmath, amsthm, mathtools}
\usepackage{graphicx}
\usepackage{systeme}
\usepackage{pgf,tikz,pgfplots}
\pgfplotsset{compat=1.15}
\usepgfplotslibrary{fillbetween}
\usepackage{mathrsfs}
\usetikzlibrary{arrows}
\usetikzlibrary{calc}
\usepackage{moresize}
\usepackage{wasysym}
\usepackage{multicol}
\usepackage{enumitem}
\renewcommand{\thechoice}{{\Large$\Circle$ \ }}
\renewcommand{\choicelabel}{\thechoice}
\newcommand{\versionmark}{$\heartsuit$}
\renewcommand{\thepartno}{\roman{partno}}
\pagestyle{headandfoot}
\setlength\parindent{0pt}
\usepackage{array}
\usepackage{amssymb} % Required for the \Box symbol
\checkboxchar{$\Box$}

\firstpageheader{CS 2050 - Exam 4 Ryland Grace\\ Spring 2026}{}{Name: \underline{\hspace{2.5in}}\\ GTID: \underline{\hspace{2.5in}}}


\runningheader{CS 2050 - Exam 4 Ryland Grace}{Page \thepage\ of \numpages}{Name: \underline{\hspace{2 in}}}
\runningheadrule

\firstpagefooter{}{}{}
\runningfooter{}{}{}
\bracketedpoints \pointsinmargin

\begin{document}

\begin{center}
\fbox{\fbox{\parbox{6in}{\centering
No notes, calculators, or other aids are allowed.  Read all directions carefully and write your answers in the space provided.\\
Taking this exam signifies you are aware of and in accordance with the Academic Honor Code of Georgia.\\ Do not separate any pages from the rest of your exam.
}}}
\end{center}
\vspace{\baselineskip}
\begin{center}
    \vspace{30mm}
    \HUGE{Version Ryland Grace}


    \numpoints\ points
\end{center}
\newpage
\vspace{\baselineskip}

\newpage

\textbf{For the following questions, please completely fill in the bubble corresponding to your chosen answer.}

\begin{questions} %------------------------------------------

% Correct answer: False; a shift cipher has only 26 possible keys (or 25 nontrivial keys).

\question[3] (new) Because a shift cipher has many possible keys, trying every possible key is not a practical method of breaking the cipher.

\begin{multicols}{2}
    \begin{choices}
        \choice True
        \choice False
    \end{choices}
\end{multicols}


\vspace{\stretch{1}}




% Correct answer: True

\question[3] (new) For every integers $n\geq 0$ and $0\leq r\leq n$, it holds
\[
C(n,r)\leq 2^n
\]

\begin{multicols}{2}
    \begin{choices}
        \choice True
        \choice False
    \end{choices}
\end{multicols}

\vspace{\stretch{1}}

% Correct answer: False; the two counts are C(8,2)=28 and C(8,3)=56.

\question[3] (new) The number of length-eight binary strings containing exactly two $1$s is equal to the number of three-element subsets of an eight-element set.

\begin{multicols}{2}
    \begin{choices}
        \choice True
        \choice False
    \end{choices}
\end{multicols}

\vspace{\stretch{1}}


% Correct answer: True; by the sum rule, there are $n_1+n_2+n_3$ choices.

\question[3] (new) A student must choose exactly one elective from three categories. The first category contains $n_1$ courses, the second contains $n_2$ courses, and the third contains $n_3$ courses. The total number of possible choices is $n_1+n_2+n_3$.

\begin{multicols}{2}
    \begin{choices}
        \choice True
        \choice False
    \end{choices}
\end{multicols}

\vspace{\stretch{1}}

% Correct answer: True; complementing every bit is a bijection between the two sets of strings.

\question[3] (new) For $0\leq r\leq n$, the number of length-$n$ binary strings containing exactly $r$ ones equals the number containing exactly $r$ zeros.

\begin{multicols}{2}
    \begin{choices}
        \choice True
        \choice False
    \end{choices}
\end{multicols}

\vspace{\stretch{1}}

\newpage

% Correct answer: B; RSA uses the public key for encryption and the private key for decryption.

\vspace{\stretch{1}}


% Correct answer: D; knowing $n$ and $e$ is insufficient to compute the private key without factoring $n$.

\question[5] (new) Which of the following statements about RSA is correct?

\begin{choices}
    \choice The security of RSA relies on keeping the modulus $n$ secret.
    \choice Anyone who knows the public key can efficiently compute the private key.
    \choice RSA is secure because multiplying large prime numbers is computationally difficult.
    \choice Factoring the modulus $n$ would allow an attacker to compute the private key.
    \choice None of the above.
\end{choices}

% Correct answer: A (5); by the Generalized Pigeonhole Principle, some category contains at least ceiling(37/8)=5 measurements.

\question[5] (new) The \textit{Hail Mary} records 37 diagnostic measurements, each assigned to exactly one of 8 categories. What is the largest number of measurements guaranteed to belong to the same category?

\begin{choices}
    \choice 5
    \choice 4
    \choice 6
    \choice 8
    \choice None of the above.
\end{choices}

\vspace{\stretch{1}}

% Correct answer: E (C(14,3)); reserve 4 packets for the command module and 1 for each other module, then distribute the remaining 11 packets.

\question[5] (new) Grace distributes 18 identical emergency food packets among four distinct ship modules. The command module must receive at least 4 packets, and each other module must receive at least 1 packet. How many distributions are possible?

\begin{choices}
    \choice $C(21,3)$
    \choice $C(17,3)$
    \choice $4^{18}$
    \choice None of the above.
    \choice $C(14,3)$
\end{choices}

\vspace{\stretch{1}}

% Correct answer: B (120); use the Product Rule: 3(2)(4)(5)=120.

\question[5] (new) A \textit{Hail Mary} diagnostic procedure has four stages. The stages can be completed in 3, 2, 4, and 5 ways, respectively, and the choice at one stage does not restrict the choices at the others. How many complete procedures are possible?

\begin{choices}
    \choice 14
    \choice 120
    \choice 60
    \choice 625
    \choice None of the above.
\end{choices}


\vspace{\stretch{1}}

% Correct answer: C (180); choose two astrophysics reports, one medical report, and two engineering reports: C(4,2)C(3,1)C(5,2)=180.

\question[5] (new) Grace selects 5 reports for a mission archive from 4 astrophysics reports, 5 engineering reports, and 3 medical reports. The archive must contain exactly 2 astrophysics reports and exactly 1 medical report. How many selections are possible?

\begin{choices}
    \choice 60
    \choice 120
    \choice 180
    \choice 360
    \choice None of the above.
\end{choices}


\vspace{\stretch{1}}

\newpage

% Correct answer: A (7!3!=30240); treat the three specified canisters as one block and arrange that block internally.

\question[5] (new) Nine distinct sample canisters are arranged in a row aboard the \textit{Hail Mary}. Three specified canisters must appear consecutively, although they may appear in any order within their block. How many arrangements are possible?

\begin{choices}
    \choice $7!3!$
    \choice $7!$
    \choice $9!/3!$
    \choice $9!-3!$
    \choice None of the above.
\end{choices}

\vspace{\stretch{1}}


% Correct answer: C (180); treat the vowels I,E,E as one block, arrange the block and four consonants in 5!/2! ways, and arrange the vowels in 3!/2! ways.

\question[5] (new) How many distinct arrangements of the word SCIENCE have all three vowels consecutive? The consonants are not required to be consecutive.

\begin{choices}
    \choice $60$
    \choice $120$
    \choice $180$
    \choice $360$
    \choice None of the above.
\end{choices}

\vspace{\stretch{1}}


% Correct answer: E (24); fix Grace's position, place Rocky in the unique opposite seat, and arrange the remaining four crew members in 4! ways.

\question[5] (new) Grace, Rocky, and four other distinct crew members sit at a circular table with six evenly spaced seats. In how many rotationally distinct seatings do Grace and Rocky sit directly opposite each other?

\begin{choices}
    \choice 12
    \choice 48
    \choice 120
    \choice None of the above.
    \choice 24
\end{choices}

\vspace{\stretch{1}}

% Correct answer: B (19); after giving each module 2 cells, count nonnegative solutions to y1+y2+y3=6 with yi at most 4: C(8,2)-3C(3,2)=19.

\question[5] (new) Grace distributes 12 identical backup fuel cells among three distinct ship modules. Each module must receive at least 2 cells and at most 6 cells. How many distributions are possible?

\begin{choices}
    \choice 16
    \choice 19
    \choice 21
    \choice 28
    \choice None of the above.
\end{choices}

\vspace{\stretch{1}}

\newpage
\textbf{Instructions:} For the following question(s), you must show \textbf{all} of your work using methods covered in lecture. Brute-force solutions will receive \textbf{no credit}. The only exception is that you may use the brute-force method to find multiplicative inverses. All other parts of the problem must be solved using techniques taught in class.

% Correct answer: $x \equiv 52 \pmod{77}$.

\question[13] Solve the following system of congruences. Give the smallest nonnegative solution.

\[
x \equiv 3 \pmod{7}
\]

\[
x \equiv 8 \pmod{11}
\]

Show all work using the Chinese Remainder Theorem.


\newpage
% Correct answer: There are 20 channels. Since 101=5(20)+1, the Generalized Pigeonhole Principle guarantees that some channel receives at least 6 transmissions.

\question[13] (new) During a communications test, the \textit{Hail Mary} sends 101 transmissions through 20 available channels. Prove that at least one channel carries at least 6 transmissions.

\newpage

\question[14] The \textit{Hail Mary} carries 16 distinct scientific samples. Give a combinatorial proof of the identity
\[
\sum_{k=4}^{12}\binom{16}{k}
=
2^{16}-2\sum_{k=0}^{3}\binom{16}{k}.
\]

In your proof:
\begin{enumerate}[label=\arabic*.]
    \item Explain what the left-hand side counts and how it counts those objects.
    \item Explain what the right-hand side counts and how it counts those objects.
    \item Explain why the two sides count the same collection of objects.
\end{enumerate}


\newpage

% change pts as needed
\question \textbf{EXTRA CREDIT:}
Design a mission patch for the \textit{Hail Mary} that includes both Earth and Rocky's home star, 40 Eridani. Submissions will be judged on a scale of 1--3. You may not raise your hand to ask a TA anything concerning this question.

\newpage

\textbf{Scratch Paper: If using this page to write a solution, clearly indicate on both this page and the original page that the solution is here.}

% =====================================================================
% INSTRUCTOR SOLUTIONS -- These comments do not appear in the PDF.
% =====================================================================
%
% Question 1: False.
% A shift cipher has 26 possible shifts, including the identity shift, so an
% exhaustive key search is practical. Even if the identity is excluded, only
% 25 nontrivial keys remain.
%
% Question 2: True.
% C(n,r) counts the r-element subsets of an n-element set. The quantity 2^n
% counts all subsets of that set. Since the r-element subsets form one part of
% the collection of all subsets, C(n,r) <= 2^n.
%
% Question 3: False.
% A length-eight binary string with exactly two 1s is uniquely determined by
% the two positions occupied by its 1s, so there are C(8,2)=28 such strings.
% There are C(8,3)=56 three-element subsets of an eight-element set. Since
% 28 is not equal to 56, the statement is false.
%
% Question 4: True.
% The three elective categories are disjoint alternatives, and the student
% chooses exactly one course. By the Sum Rule, the total is n_1+n_2+n_3.
%
% Question 5: True.
% Complement every bit in a string. This operation is a bijection from strings
% with exactly r ones to strings with exactly r zeros, and applying it twice
% returns the original string.
%
% Question 6: Correct choice D: factoring the modulus n would allow an attacker
% to compute the private key.
% If n=pq is factored, the attacker can compute phi(n)=(p-1)(q-1) and then find
% the private exponent d satisfying ed congruent to 1 modulo phi(n). The modulus
% n is part of the public key, so it is not secret. Efficiently deriving the
% private key from the public key is precisely what RSA is intended to prevent.
% Multiplication of large primes is efficient; reversing it by factoring n is
% the computationally difficult operation.
%
% Question 7: Correct choice A: 5.
% Place the 37 measurements into 8 category boxes. The Generalized Pigeonhole
% Principle guarantees at least ceiling(37/8)=5 measurements in one category.
% The bound is sharp because counts 5,5,5,5,5,4,4,4 total 37 and no category
% contains more than 5 measurements.
%
% Question 8: Correct choice E: C(14,3)=364.
% Reserve 4 packets for the command module and 1 packet for each of the other
% three modules. This uses 7 packets and leaves 11. If y_i denotes additional
% packets, then y_1+y_2+y_3+y_4=11 with every y_i nonnegative. Stars and bars
% gives C(11+4-1,4-1)=C(14,3)=364.
%
% Question 9: Correct choice B: 120.
% The four independent stages have 3, 2, 4, and 5 choices. By the Product Rule,
% there are 3(2)(4)(5)=120 complete procedures.
%
% Question 10: Correct choice C: 180.
% Choose 2 of the 4 astrophysics reports, 1 of the 3 medical reports, and the
% remaining 2 reports from the 5 engineering reports. The Product Rule gives
% C(4,2)C(3,1)C(5,2)=6(3)(10)=180.
%
% Question 11: Correct choice A: 7!3!=30,240.
% Treat the three specified canisters as one block. The block and the other six
% canisters form seven distinct objects, which can be ordered in 7! ways. The
% three canisters inside the block can be ordered in 3! ways. Multiply.
%
% Question 12: Correct choice C: 180.
% SCIENCE has vowels I,E,E and consonants S,C,N,C. Treat the three vowels as one
% block. The vowel block and four consonants form five objects, with the letter
% C repeated, so they have 5!/2!=60 arrangements. The vowels have 3!/2!=3
% internal arrangements. Thus the total is (60)(3)=180.
%
% Question 13: Correct choice E: 24.
% Fix Grace's seat to remove rotational symmetry. Rocky must occupy the unique
% seat opposite Grace. The other four distinct crew members may then be placed
% in the remaining seats in 4!=24 ways.
%
% Question 14: Correct choice B: 19.
% First give 2 cells to every module. For the remaining 6 cells, solve
% y_1+y_2+y_3=6 with 0<=y_i<=4. Without the upper bounds there are C(8,2)=28
% solutions. If a particular y_i>=5, set z_i=y_i-5; the new sum is 1 and has
% C(3,2)=3 solutions. There are three choices of the violating variable, and no
% two variables can both be at least 5. Hence 28-3C(3,2)=28-9=19.
%
% Question 15: Correct answer: x congruent to 52 modulo 77; the smallest
% nonnegative solution is 52.
% Write x=3+7t. The second congruence gives 3+7t congruent to 8 modulo 11, so
% 7t congruent to 5 modulo 11. Since 8 is the inverse of 7 modulo 11, multiply
% by 8 to obtain t congruent to 40 congruent to 7 modulo 11. Therefore
% x=3+7(7)=52. All solutions are x congruent to 52 modulo 7(11)=77.
%
% Question 16:
% Regard the 20 channels as boxes and the 101 transmissions as objects. If every
% channel carried at most 5 transmissions, at most 20(5)=100 transmissions could
% be sent. This contradicts the total of 101. Therefore some channel carries at
% least ceiling(101/20)=6 transmissions.
%
% Question 17:
% Both sides count collections containing at least 4 and at most 12 of the 16
% distinct samples. The left side counts directly by collection size. For each
% k from 4 through 12, there are C(16,k) collections, and the Sum Rule gives the
% left-hand side.
% There are 2^16 collections of samples in total. The invalid collections have
% sizes 0,1,2,3,13,14,15, or 16. The collections of sizes 0 through 3 contribute
% sum from k=0 to 3 of C(16,k). Taking complements is a bijection from a
% k-element collection to a (16-k)-element collection, so the collections of
% sizes 13 through 16 contribute the same number. Subtracting both groups of
% invalid collections gives the right-hand side. Both methods count exactly the
% same valid collections, proving the identity. Their common value is
% 2^16-2[C(16,0)+C(16,1)+C(16,2)+C(16,3)]
% =65536-2(1+16+120+560)=64142.
%
% Question 18:
% Answers vary. Award 1--3 points according to the announced art rubric if the
% patch includes Earth and 40 Eridani.
% =====================================================================

\end{questions}

\end{document}
